\documentclass[10pt, aspectratio=169]{beamer}
\usetheme{Madrid}
\usecolortheme{default}

% Packages
\usepackage[utf8]{inputenc}
\usepackage[T1]{fontenc}
\usepackage{lmodern}
\usepackage{graphicx}
\usepackage{booktabs}
\usepackage{amsmath}

% Metadata
\title[Trajectory Gen. for Redundant Welding Robot]{Trajectory Generation for Redundant Welding Robot with Positioner}
\author{Xuan Dinh Nguyen}
\institute[CTU in Prague]{
    Czech Technical University in Prague\\
    Faculty of Electrical Engineering\\
    Department of Cybernetics
}
\date{\today}

\begin{document}

% Slide 1: Title
\begin{frame}
    \titlepage
\end{frame}

% Slide 2: Motivation
\begin{frame}{Motivation \& Problem Statement}
    \begin{itemize}
        \item \textbf{Complex Geometry:} Robotic welding often involves complex parts where a standard 6-DOF arm cannot reach all seams due to collisions or kinematic limits.
        \item \textbf{Redundancy:} Adding a positioner (e.g., 2-DOF) increases the total degrees of freedom (9-DOF system), creating kinematic redundancy.
        \item \textbf{Goal:} Develop a complete software solution to automatically plan motions and generate trajectories for this 9-DOF welding cell directly from standard G-code.
    \end{itemize}
\end{frame}

% Slide 3: System Overview
\begin{frame}{System Overview}
    \begin{columns}
        \column{0.5\textwidth}
        \textbf{Hardware Setup:}
        \begin{itemize}
            \item 6-DOF Industrial Robotic Arm
            \item 2-DOF Rotary Positioner
            \item 1 External Rotational Joint
        \end{itemize}
        
        \vspace{0.5cm}
        \textbf{Software Stack:}
        \begin{itemize}
            \item ROS 2 (Humble)
            \item MoveIt 2
            \item Ceres Solver for optimisation
        \end{itemize}
        
        \column{0.5\textwidth}
        \textbf{Workflow:}
        \begin{enumerate}
            \item Input: Standard G-code
            \item Inverse Kinematics (Analytical)
            \item Redundancy Resolution (Optimisation)
            \item Collision-free repositioning
            \item Output: Time-parametrised joint trajectory
        \end{enumerate}
    \end{columns}
\end{frame}

% Slide 4: Kinematics
\begin{frame}{Kinematic Approach}
    \begin{itemize}
        \item \textbf{Analytical Foundation:}
        \begin{itemize}
            \item Closed-form inverse kinematics (IK) derived for both the arm and the positioner.
            \item Faster and more reliable than purely numerical IK solvers.
        \end{itemize}
        \item \textbf{Handling Redundancy:}
        \begin{itemize}
            \item The system has more DOFs than strictly necessary for a 6D pose.
            \item We use non-linear optimisation (Ceres Solver) to find the best configuration among infinite possibilities.
        \end{itemize}
    \end{itemize}
\end{frame}

% Slide 5: The Key Concept - Pose Relaxation
\begin{frame}{Pose Relaxation Strategy}
    To navigate singularities and restrictive geometries, the strict rigid welding pose is relaxed in 5 DOFs:
    
    \vspace{0.3cm}
    \begin{columns}
        \column{0.5\textwidth}
        \textbf{Positioner Relaxation:}
        \begin{itemize}
            \item Roll ($\xi_{p}$)
            \item Slope ($\theta_{p}$)
        \end{itemize}
        
        \column{0.5\textwidth}
        \textbf{Arm Torch Relaxation:}
        \begin{itemize}
            \item Work angle ($\alpha_{\mathrm{arm}}$)
            \item Travel angle ($\beta_{\mathrm{arm}}$)
            \item Torch spin ($\gamma_{\mathrm{arm}}$)
        \end{itemize}
    \end{columns}
    
    \vspace{0.5cm}
    \begin{center}
        \textit{This gives the optimiser the freedom to smoothly adapt the tool orientation without violating process requirements.}
    \end{center}
\end{frame}

% Slide 6: Trajectory Optimisation
\begin{frame}{Trajectory Optimisation}
    The optimal configuration is found by minimising an objective function that balances multiple criteria:
    \begin{itemize}
        \item \textbf{Manipulability:} Keeping the robot far from singular configurations.
        \item \textbf{Joint Limits:} Penalising configurations near the mechanical bounds.
        \item \textbf{Collision Avoidance:} Maximising clearance from obstacles.
        \item \textbf{Smoothness:} Ensuring smooth transitions between waypoints across the continuous weld seam.
    \end{itemize}
\end{frame}

% Slide 7: Repositioning (Free Space)
\begin{frame}{Rapid Repositioning}
    \begin{itemize}
        \item Between welding seams (air moves), the robot must reposition quickly and safely.
        \item \textbf{Method:} RRT-Connect planner.
        \item \textbf{Advantages:}
        \begin{itemize}
            \item Sampling-based approach perfectly suited for high-dimensional configuration spaces (9-DOF).
            \item Guarantees collision-free paths around complex workpieces.
        \end{itemize}
    \end{itemize}
\end{frame}

% Slide 8: Implementation \& Results
\begin{frame}{Implementation \& Results}
    \begin{itemize}
        \item \textbf{Modularity:} Built as a set of modular ROS 2 nodes, integrating seamlessly with MoveIt 2.
        \item \textbf{User Interface:} Custom GUI (PyQt5) for easy operation and monitoring.
        \item \textbf{Validation:} Successfully generates smooth, collision-free trajectories for complex test geometries that would fail with naive 6-DOF approaches.
    \end{itemize}
\end{frame}

% Slide 9: Conclusion
\begin{frame}{Conclusion}
    \begin{itemize}
        \item Developed a robust, automated pipeline for a 9-DOF redundant welding cell.
        \item Analytical kinematics combined with non-linear optimisation provides both reliability and flexibility.
        \item The 5-DOF relaxation strategy is the key enabler for welding complex parts.
        \item Fully integrated into a modern ROS 2 software stack.
    \end{itemize}
    
    \vspace{0.5cm}
    \begin{center}
        \Large \textbf{Thank you for your attention!}
    \end{center}
\end{frame}

\end{document}
